% Cayley's Collected Mathematical Papers, Vol. II, Pages 141--184 (PDF 157--200)
% Transcribed from the original 1889 Cambridge University Press edition
% Covers: tail of paper 128, papers 129--133 (partial 133).
\documentclass[11pt,a4paper]{article}
\usepackage[utf8]{inputenc}
\usepackage[T1]{fontenc}
\usepackage[french,english]{babel}
\usepackage{amsmath,amssymb,amsthm}
\usepackage{geometry}
\usepackage{graphicx}
\usepackage{fancyhdr}
\usepackage[protrusion=true,expansion=false]{microtype}
\usepackage{hyperref}

\geometry{margin=1in}

\pagestyle{fancy}
\fancyhf{}
\fancyhead[C]{\thepage}
\renewcommand{\headrulewidth}{0.4pt}

\providecommand{\md}{\,\mathrm{mod.}\,}
\providecommand{\mat}[1]{\mathbf{#1}}

\title{Cayley's Collected Mathematical Papers, Vol.~II, pp.\ 141--184}
\author{Reconstruction with OCR skeleton and PNG verification}
\date{}

\emergencystretch=3em
\tolerance=600
\begin{document}
\maketitle

%==============================================================================
\section*{128.\ (continued)}
\addcontentsline{toc}{section}{128. Developments on the Porism of the In-and-Circumscribed Polygon (continued)}
\subsection*{Developments on the Porism of the In-and-Circumscribed Polygon (continued).}

[From the \textit{Philosophical Magazine}, vol.\ vii.\ (1854), pp.\ 339--345.]

\medskip

\noindent (Continuation from p.~140.)

\medskip

which, he remarks, is satisfied by $r = -p$ and $r = -\dfrac{pq}{p+q}$, and that consequently the rationalized equation will divide by $p+r$ and $pq - r(p+q)$; and he finds, after the division,
\[
p^2 q^2 + p^2 q^2 (p+q)\, r - pq (p+q)^2\, r^2 - (p+q)(p-q)^2\, r^3 = 0,
\]
which, restoring for $p$, $q$ their values $R+a$, $R-a$, is the very equation above found.

\medskip

The form given by Steiner (\textit{Crelle}, t.\ ii.\ p.\ 289) is
\[
r(R-a)^2 = (R+a)\sqrt{(R-r+a)(R-r-a)} + (R+a)\sqrt{(R-r-a)\cdot 2R},
\]
which, putting $p$, $q$ instead of $R+a$, $R-a$, is
\[
qr = p\sqrt{(p-r)(q-r)} + p\sqrt{(q-r)(q+p)};
\]
and Jacobi has shown in his memoir, ``Anwendung der elliptischen Transcendenten u.\ s.\ w.,'' \textit{Crelle}, t.\ iii.\ [1828] p.\ 376, that the rationalized equation divides (like that of Fuss) by the factor $pq - (p+q)r$, and becomes by that means identical with the rational equation given by Fuss.

\medskip

In the case of two concentric circles $a = 0$, and putting $\dfrac{r}{R} = M$, we have
\[
\sqrt{1-M} = \sqrt{1-M}.
\]
This is, in fact, the very formula which corresponds to the general case of two conics having double contact. For suppose that the polygon is inscribed in the conic $U = 0$, and circumscribed about the conic $U + P^2 = 0$, we have then to find the discriminant of $\xi U + U + P^2$, i.e.\ of $(1+\xi)U + P^2$. Let $K$ be the discriminant of $U$, and let $F$ be what the polar reciprocal of $U$ becomes when the variables are replaced by the coefficients of $P$, or, what is the same thing, let $-F$ be the determinant obtained by bordering $K$ (considered as a matrix) with the coefficients of $P$. The discriminant of $(1+\xi)U + P^2$ is $(1+\xi)^3 K + (1+\xi)^2 F$, i.e.\ it is
\[
(1+\xi)^2 K \bigl[ (1+\xi) + \tfrac{F}{K}\bigr],\quad\text{i.e.\ } (1+\xi)^2 K (1 + M\xi),
\]
where $M = 1 + \dfrac{F}{K}$; or, what is the same thing, $M$ is the discriminant of $U$ divided by the discriminant of $U + P^2$. And $M$ having this meaning, the condition of there being inscribed in the conic $U$ an infinity of $n$-gons circumscribed about the conic $U + P^2 = 0$, is found by means of the series
\[
A + B\xi + C\xi^2 + D\xi^3 + E\xi^4 + \&\mathrm{c}. = (1+\xi)\sqrt{1+M\xi}.
\]

\medskip

We have, therefore,
\begin{align*}
A &= 1,\\
2B &= M + 2,\\
8C &= -M^2 + 4M,\\
16 D &= M^3 - 2M^2,\\
-128 E &= 5M^4 - 8M^3,\\
&\,\&\mathrm{c}.,\\
1024\,(CE - D^2) &= M^4(M^2 - 12M + 16),\\
&\,\&\mathrm{c}.
\end{align*}

\medskip

Hence for the triangle, quadrangle and pentagon, the conditions are
\begin{description}
\item[I.] For the triangle,
\[ M + 2 = 0. \]
\item[II.] For the quadrangle,
\[ M - 4 = 0. \]
\item[III.] For the pentagon,
\[ M^2 - 12M + 16 = 0; \]
\end{description}
and so on.

\medskip

It is worth noticing, that, in the case of two conics having a 4-point contact, we have $F = 0$, and consequently $M = 1$. The discriminant is therefore $(1 + \xi)^3$, and as this does not contain any variable parameter, the conics cannot be determined so that there may be for a given value of $n$ (nor, indeed, for any value whatever of $n$) an infinity of $n$-gons inscribed in the one conic, and circumscribed about the other conic.

\medskip

The geometrical properties of a triangle, \&c.\ inscribed in a conic and circumscribed about another conic, these two conics having double contact with each other, are at once obtained from those of the system in which the two conics are replaced by concentric circles. Thus, in the case of a triangle, if $ABC$ be the triangle, and $\alpha$, $\beta$, $\gamma$ be the points of contact of the sides with the inscribed conic, then the tangents to the circumscribed conic at $A$, $B$, $C$ meet the opposite sides $BC$, $CA$, $AB$ in points lying in the chord of contact, the lines $A\alpha$, $B\beta$, $C\gamma$ meet in the pole of contact, and so on.

\medskip

In the case of a quadrangle, if $ACEG$ be the quadrangle, and $b$, $d$, $f$, $h$ the points of contact with the inscribed conic, then the tangents to the circumscribed conic at the pair of opposite angles $A$, $E$ and the corresponding diagonal $CG$, and in like manner the tangents at the pair of opposite angles $C$, $G$ and the corresponding diagonal $AE$, meet in the chord of contact. Again, the pairs of opposite sides $AC$, $EG$, and the line $dh$ joining the points of contact of the other two sides with the inscribed conic, and the pairs of opposite sides $AG$, $CE$, and the line $hf$ joining the points of contact of the other two sides with the inscribed conic, meet in the chord of contact. The diagonals $AE$, $CG$, and the lines $bf$, $dh$ through the points of contact of pairs of opposite sides with the inscribed conic, meet in the pole of contact, \&c.

\medskip

The beautiful systems of `focal relations' for regular polygons (in particular for the pentagon and the hexagon), given in Sir W.\ R.\ Hamilton's \textit{Lectures on Quaternions}, [Dublin, 1853] Nos.\ 379--393, belong, it is clear, to polygons which are inscribed in and circumscribed about two conics having double contact with each other. In fact, the focus of a conic is a point such that the lines joining such point with the circular points at infinity (i.e.\ the points in which a circle is intersected by the line infinity) are tangents to the conic. In the case of two concentric circles, these are to be considered as touching in the circular points at infinity; and consequently, when the concentric circles are replaced by two conics having double contact, the circular points at infinity are replaced by the points of contact of the two conics.

\medskip

Thus, in the figure (which is simply Sir W.\ R.\ Hamilton's figure 81 put into perspective), the system of relations
\begin{align*}
F, G(\,..)\, &ABCI,\\
G, H(\,..)\, &BCDK,\\
H, I(\,..)\, &CDEF,\\
I, K(\,..)\, &DEAG,\\
K, F(\,..)\, &EABH,
\end{align*}
will mean, $F, G(\,..)\,ABCI$, that there is a conic inscribed in the quadrilateral $ABCI$ such that the tangents to this conic through the points $F$ and $G$ pass two and two through the points of contact of the circumscribed and the inscribed conics, and similarly for the other relations of the system. As the figure is drawn, the tangents in question are of course (as the tangents through the foci in the case of the two concentric circles) imaginary.

\medskip
\noindent\textit{2 Stone Buildings, March 7, 1854.}

%==============================================================================
\section*{129.}
\addcontentsline{toc}{section}{129. On the Porism of the In-and-Circumscribed Triangle, and on an Irrational Transformation of Two Ternary Quadratic Forms Each into Itself}
\subsection*{On the Porism of the In-and-Circumscribed Triangle, and on an Irrational \\ Transformation of Two Ternary Quadratic Forms each into itself.}

[From the \textit{Philosophical Magazine}, vol.\ ix.\ (1855), pp.\ 513--517.]

\medskip

There is an irrational transformation of two ternary quadratic forms each into itself, based upon the solution of the following geometrical problem,

\medskip

\textit{Given that the line}
\[ lx + my + nz = 0 \]
\textit{meets the conic}
\[ (a, b, c, f, g, h \mid x, y, z)^2 = 0 \]
\textit{in the point $(x_1, y_1, z_1)$; to find the other point of intersection.}

\medskip

The solution is exceedingly simple. Take $(x_2, y_2, z_2)$ for the coordinates of the other point of intersection, we must have identically with respect to $x$, $y$, $z$,
\begin{align*}
(a, \ldots \mid x, y, z)^2 \cdot (\mathfrak{A}, \ldots \mid l, m, n)^2 - k(lx + my + nz)^2 \\
\quad = (a, \ldots \mid x_1, y_1, z_1 \mid x, y, z)\cdot(a, \ldots \mid x_2, y_2, z_2 \mid x, y, z)
\end{align*}
to a constant factor pr\`es.

\medskip

Assume successively $x, y, z = \mathfrak{A}, \mathfrak{H}, \mathfrak{G}$; $\mathfrak{H}, \mathfrak{B}, \mathfrak{F}$; $\mathfrak{G}, \mathfrak{F}, \mathfrak{C}$; it follows that
\begin{align*}
x_2 : y_2 : z_2 = \;& y_1 z_1 \bigl[ \mathfrak{A}(a, \ldots \mid l, m, n)^2 - (\mathfrak{A}l + \mathfrak{H}m + \mathfrak{G}n)^2 \bigr]\\
:\;& z_1 x_1 \bigl[ \mathfrak{B}(a, \ldots \mid l, m, n)^2 - (\mathfrak{H}l + \mathfrak{B}m + \mathfrak{F}n)^2 \bigr]\\
:\;& x_1 y_1 \bigl[ \mathfrak{C}(a, \ldots \mid l, m, n)^2 - (\mathfrak{G}l + \mathfrak{F}m + \mathfrak{C}n)^2 \bigr];
\end{align*}
or, what is the same thing,
\begin{align*}
x_2 : y_2 : z_2 = \;& y_1 z_1 (bn^2 + cm^2 - 2fmn)\\
:\;& z_1 x_1 (cl^2 + an^2 - 2gnl)\\
:\;& x_1 y_1 (am^2 + bl^2 - 2hlm).
\end{align*}

\medskip

It is not necessary for the present purpose, but it may be as well to give the corresponding solution of the problem

\medskip

\textit{Given that one of the tangents through the point $(\xi, \eta, \zeta)$ to the conic}
\[ (a, b, c, f, g, h \mid x, y, z)^2 = 0 \]
\textit{is the line $l_1 x + m_1 y + n_1 z = 0$; to find the equation to the other tangent.}

\medskip

Let $l_2 x + m_2 y + n_2 z = 0$ be the other tangent, then
\begin{align*}
(a, \ldots \mid \xi, \eta, \zeta)^2 \cdot (a, \ldots \mid x, y, z)^2 - \bigl[(a, \ldots \mid \xi, \eta, \zeta \mid x, y, z)\bigr]^2\\
= (l_1 x + m_1 y + n_1 z)(l_2 x + m_2 y + n_2 z)
\end{align*}
to a constant factor pr\`es. Assume successively $y = 0, z = 0$; $z = 0, x = 0$; $x = 0, y = 0$; then we have
\begin{align*}
l_2 : m_2 : n_2 = \;& m_1 n_1 \bigl[ a (a, \ldots \mid \xi, \eta, \zeta)^2 - (a\xi + h\eta + g\zeta)^2 \bigr]\\
:\;& n_1 l_1 \bigl[ b (a, \ldots \mid \xi, \eta, \zeta)^2 - (h\xi + b\eta + f\zeta)^2 \bigr]\\
:\;& l_1 m_1 \bigl[ c (a, \ldots \mid \xi, \eta, \zeta)^2 - (g\xi + f\eta + c\zeta)^2 \bigr];
\end{align*}
or, as they may be more simply written,
\begin{align*}
l_2 : m_2 : n_2 = \;& m_1 n_1 (\mathfrak{B}\zeta^2 + \mathfrak{C}\eta^2 - 2\mathfrak{F}\eta\zeta)\\
:\;& n_1 l_1 (\mathfrak{C}\xi^2 + \mathfrak{A}\zeta^2 - 2\mathfrak{G}\zeta\xi)\\
:\;& l_1 m_1 (\mathfrak{A}\eta^2 + \mathfrak{B}\xi^2 - 2\mathfrak{H}\xi\eta).
\end{align*}

\medskip

Returning now to the solution of the first problem, I shall for the sake of simplicity consider the formul\ae{} obtained by taking for the equation of the conic,
\[ \alpha x^2 + \beta y^2 + \gamma z^2 = 0. \]

\medskip

We see, therefore, that if this conic be intersected by the line $lx + my + nz = 0$ in the points $(x_1, y_1, z_1)$ and $(x_2, y_2, z_2)$, then
\begin{align*}
x_2 : y_2 : z_2 = \;& y_1 z_1 (\gamma m^2 + \beta n^2)\\
:\;& z_1 x_1 (\alpha n^2 + \gamma l^2)\\
:\;& x_1 y_1 (\beta l^2 + \alpha m^2).
\end{align*}

\medskip

We have, in fact, identically
\begin{align*}
& l y_1 z_1 (\beta n^2 + \gamma m^2) + m z_1 x_1 (\gamma l^2 + \alpha n^2) + n x_1 y_1 (\alpha m^2 + \beta l^2)\\
&\quad = (\alpha m n x_1 + \beta n l y_1 + \gamma l m z_1)(l x_1 + m y_1 + n z_1) - lmn (\alpha x_1^2 + \beta y_1^2 + \gamma z_1^2).
\end{align*}
Also,
\begin{align*}
& \alpha y_1^2 z_1^2 (\beta n^2 + \gamma m^2)^2 + \beta z_1^2 x_1^2 (\gamma l^2 + \alpha n^2)^2 + \gamma x_1^2 y_1^2 (\alpha m^2 + \beta l^2)^2\\
&\quad = \alpha\beta\gamma\bigl[ - l^2 x_1^2 - m^2 y_1^2 - n^2 z_1^2 \\
&\qquad + (m y_1 + n z_1)^2 l^2 x_1^2 + (n z_1 + l x_1)^2 m^2 y_1^2 + (l x_1 + m y_1)^2 n^2 z_1^2 - 2 lmn x_1 y_1 z_1 (l x_1 + m y_1 + n z_1)\bigr]\\
&\qquad - (l^2 \beta\gamma x_1^2 + m^2 \gamma\alpha y_1^2 + n^2 \alpha\beta z_1^2)(\alpha x_1^2 + \beta y_1^2 + \gamma z_1^2);
\end{align*}
which show that if $l x_1 + m y_1 + n z_1 = 0$ and $\alpha x_1^2 + \beta y_1^2 + \gamma z_1^2 = 0$, then also $l x_2 + m y_2 + n z_2 = 0$ and $\alpha x_2^2 + \beta y_2^2 + \gamma z_2^2 = 0$: this is, of course, as it should be.

\medskip

I shall now consider $l$, $m$, $n$ as given functions of $x_1$, $y_1$, $z_1$ satisfying identically the equations
\begin{align*}
l x_1 + m y_1 + n z_1 &= 0,\\
l^2 bc + m^2 ca + n^2 ab &= 0,
\end{align*}
equations which express that $lx + my + nz = 0$ is the tangent from the point $(x_1, y_1, z_1)$ to the conic $ax^2 + by^2 + cz^2 = 0$. And I shall take for $\alpha$, $\beta$, $\gamma$ the following values, viz.
\begin{align*}
\alpha &= ax_1^2 + by_1^2 + cz_1^2 - a(x_1^2 + y_1^2 + z_1^2),\\
\beta &= ax_1^2 + by_1^2 + cz_1^2 - b(x_1^2 + y_1^2 + z_1^2),\\
\gamma &= ax_1^2 + by_1^2 + cz_1^2 - c(x_1^2 + y_1^2 + z_1^2);
\end{align*}
so that $x_1$, $y_1$, $z_1$ continuing absolutely indeterminate, we have identically $\alpha x_1^2 + \beta y_1^2 + \gamma z_1^2 = 0$. Also taking $\Theta$ as a function of $x_1$, $y_1$, $z_1$, the value of which will be subsequently given, I write
\begin{align*}
x_2 &= \Theta y_1 z_1 (\beta n^2 + \gamma m^2),\\
y_2 &= \Theta z_1 x_1 (\gamma l^2 + \alpha n^2),\\
z_2 &= \Theta x_1 y_1 (\alpha m^2 + \beta l^2);
\end{align*}
so that $x_1$, $y_1$, $z_1$ are arbitrary, and $x_2$, $y_2$, $z_2$ are taken to be determinate functions of $x_1$, $y_1$, $z_1$. The point $(x_2, y_2, z_2)$ is geometrically connected with the point $(x_1, y_1, z_1)$ as follows, viz.\ $(x_2, y_2, z_2)$ is the point in which the tangent through $(x_1, y_1, z_1)$ to the conic $ax^2 + by^2 + cz^2 = 0$ meets the conic passing through the point $(x_1, y_1, z_1)$ and the points of intersection of the conics $ax^2 + by^2 + cz^2 = 0$ and $x^2 + y^2 + z^2 = 0$. Consequently, in the particular case in which $(x_1, y_1, z_1)$ is a point on the conic $x^2 + y^2 + z^2 = 0$, the point $(x_2, y_2, z_2)$ is the point in which this conic is met by the tangent through $(x_1, y_1, z_1)$ to the conic $ax^2 + by^2 + cz^2 = 0$.

\medskip

It has already been seen that $l x_1 + m y_1 + n z_1 = 0$ and $\alpha x_1^2 + \beta y_1^2 + \gamma z_1^2 = 0$ identically; consequently we have identically $l x_2 + m y_2 + n z_2 = 0$ and $\alpha x_2^2 + \beta y_2^2 + \gamma z_2^2 = 0$. The latter equation, written under the form
\[
(a x_1^2 + b y_1^2 + c z_1^2)(x_2^2 + y_2^2 + z_2^2) - (x_1^2 + y_1^2 + z_1^2)(a x_2^2 + b y_2^2 + c z_2^2) = 0,
\]
shows that if $x_2$, $y_2$, $z_2$ are such that $x_2^2 + y_2^2 + z_2^2 = x_1^2 + y_1^2 + z_1^2$, then that also $a x_2^2 + b y_2^2 + c z_2^2 = a x_1^2 + b y_1^2 + c z_1^2$. I proceed to determine $\Theta$ so that we may have $x_2^2 + y_2^2 + z_2^2 = x_1^2 + y_1^2 + z_1^2$. We obtain immediately
\begin{align*}
\tfrac{1}{\Theta^2}(x_2^2 + y_2^2 + z_2^2) ={}& (l^2 x_1^2 + m^2 y_1^2 + n^2 z_1^2)(\alpha^2 x_1^2 + \beta^2 y_1^2 + \gamma^2 z_1^2)\\
& - (\alpha^2 l^2 x_1^4 + \beta^2 m^2 y_1^4 + \gamma^2 n^2 z_1^4 - 2\beta\gamma m n l^2 y_1^2 z_1^2 - 2\gamma\alpha n l m^2 z_1^2 x_1^2 - 2\alpha\beta l m n^2 x_1^2 y_1^2);
\end{align*}
write for a moment
\[
ax_1^2 + by_1^2 + cz_1^2 = p,\quad x_1^2 + y_1^2 + z_1^2 = q,\quad\text{so that}\quad \alpha = p - aq,\;\beta = p - bq,\;\gamma = p - cq,
\]
then
\begin{align*}
\alpha^2 x_1^2 + \beta^2 y_1^2 + \gamma^2 z_1^2 &= q p^2 - 2 p\cdot pq + (a^2 x_1^2 + b^2 y_1^2 + c^2 z_1^2)\,q^2\\
&= q\bigl\{(a^2 x_1^2 + b^2 y_1^2 + c^2 z_1^2)\,q - p^2\bigr\}\\
&= q\bigl\{(b - c)^2 y_1^2 z_1^2 + (c - a)^2 z_1^2 x_1^2 + (a - b)^2 x_1^2 y_1^2\bigr\}.
\end{align*}
Moreover,
\begin{align*}
&\alpha^2 l^2 x_1^4 + \beta^2 m^2 y_1^4 + \gamma^2 n^2 z_1^4 - 2\beta\gamma m n l^2 y_1^2 z_1^2 - 2\gamma\alpha n l m^2 z_1^2 x_1^2 - 2\alpha\beta l m n^2 x_1^2 y_1^2\\
&\quad = p^2 \bigl\{ l^2 x_1^4 + m^2 y_1^4 + n^2 z_1^4 - 2 m n y_1^2 z_1^2 - 2 n l z_1^2 x_1^2 - 2 l m x_1^2 y_1^2\bigr\}\\
&\qquad - 2 pq \bigl\{a l^2 x_1^4 + b m^2 y_1^4 + c n^2 z_1^4 - (b + c) m n y_1^2 z_1^2 - (c + a) n l z_1^2 x_1^2 - (a + b) l m x_1^2 y_1^2\bigr\}\\
&\qquad + q^2 \bigl\{ a^2 l^2 x_1^4 + b^2 m^2 y_1^4 + c^2 n^2 z_1^4 - 2 b c m n y_1^2 z_1^2 - 2 c a n l z_1^2 x_1^2 - 2 a b l m x_1^2 y_1^2\bigr\},
\end{align*}
the first line of which vanishes in virtue of the equation $l x_1 + m y_1 + n z_1 = 0$; we have therefore
\begin{align*}
\tfrac{1}{\Theta^2}(x_2^2 + y_2^2 + z_2^2) - (x_1^2 + y_1^2 + z_1^2) ={}& (l^2 x_1^2 + m^2 y_1^2 + n^2 z_1^2)\bigl\{(b-c)^2 y_1^2 z_1^2 + (c-a)^2 z_1^2 x_1^2 + (a-b)^2 x_1^2 y_1^2\bigr\}\\
&+ 2 p q \{\cdots\} - q^2 \{\cdots\}.
\end{align*}

\medskip

Hence reducing the function on the right-hand side, and putting
\[
\tfrac{1}{\Theta^2}(x_2^2 + y_2^2 + z_2^2) \div (x_1^2 + y_1^2 + z_1^2) = 1,
\]
we have
\begin{align*}
\tfrac{1}{\Theta^2} ={}& a^2 l^4 x_1^4 + b^2 m^4 y_1^4 + c^2 n^4 z_1^4\\
&+ (c^2 m^4 - 2 b^2 m^2 n^2)\,y_1^2 z_1^2 + (a^2 n^4 - 2 c^2 n^2 l^2)\,z_1^2 x_1^2 + (b^2 l^4 - 2 a^2 l^2 m^2)\,x_1^2 y_1^2\\
&+ (b^2 n^4 - 2 c^2 m^2 n^2)\,y_1^2 z_1^2 + (c^2 l^4 - 2 a^2 n^2 l^2)\,z_1^2 x_1^2 + (a^2 m^4 - 2 b^2 l^2 m^2)\,x_1^2 y_1^2\\
&+ \bigl\{l^4 (b-c)^2 + m^4 (c-a)^2 + n^4 (a-b)^2\\
&\qquad + 2 m^2 n^2 (bc - ca - ab) + 2 n^2 l^2 (-bc + ca - ab) + 2 l^2 m^2 (-bc - ca + ab)\bigr\}\,x_1^2 y_1^2 z_1^2.
\end{align*}
The value of $\Theta$ might probably be expressed in a more simple form by means of the equations $l x_1 + m y_1 + n z_1 = 0$ and $l^2 bc + m^2 ca + n^2 ab = 0$, even without solving these equations; but this I shall not at present inquire into.

\medskip

Recapitulating, $l$, $m$, $n$ are considered as functions of $x_1$, $y_1$, $z_1$ determined (to a common factor pr\`es) by the equations
\begin{align*}
l x_1 + m y_1 + n z_1 &= 0,\\
l^2 bc + m^2 ca + n^2 ab &= 0;
\end{align*}
$\Theta$ is determined as above, and then writing
\begin{align*}
\alpha &= a x_1^2 + b y_1^2 + c z_1^2 - a(x_1^2 + y_1^2 + z_1^2),\\
\beta &= a x_1^2 + b y_1^2 + c z_1^2 - b(x_1^2 + y_1^2 + z_1^2),\\
\gamma &= a x_1^2 + b y_1^2 + c z_1^2 - c(x_1^2 + y_1^2 + z_1^2),
\end{align*}
we have
\begin{align*}
x_2 &= \Theta y_1 z_1 (\beta n^2 + \gamma m^2),\\
y_2 &= \Theta z_1 x_1 (\gamma l^2 + \alpha n^2),\\
z_2 &= \Theta x_1 y_1 (\alpha m^2 + \beta l^2);
\end{align*}
and these values give
\begin{align*}
l x_2 + m y_2 + n z_2 &= 0,\\
x_2^2 + y_2^2 + z_2^2 &= x_1^2 + y_1^2 + z_1^2,\\
a x_2^2 + b y_2^2 + c z_2^2 &= a x_1^2 + b y_1^2 + c z_1^2.
\end{align*}

\medskip

In connexion with the subject I may add the following transformation, viz.\ if
\[
3 \sqrt{\alpha}\, x' = \sqrt{3 \alpha}\,(y - z) + \sqrt{(3 \alpha - 2 \beta)(x^2 + y^2 + z^2) + 2 \beta (yz + zx + xy)},
\]
then reciprocally
\[
3 \sqrt{\beta}\, x = - \sqrt{3 \beta}\,(y' - z') + \sqrt{(3 \beta - 2 \alpha)(x'^2 + y'^2 + z'^2) + 2 \alpha (y'z' + z'x' + x'y')},
\]
and
\begin{align*}
x^2 + y^2 + z^2 &= x'^2 + y'^2 + z'^2,\\
\beta\,(x^2 + y^2 + z^2 - yz - zx - xy) &= \alpha\,(x'^2 + y'^2 + z'^2 - y'z' - z'x' - x'y').
\end{align*}
Suppose $1 + \rho + \rho^2 = 0$, then
\[
x^2 + y^2 + z^2 - yz - zx - xy = (x + \rho y + \rho^2 z)(x + \rho^2 y + \rho z);
\]
and in fact
\begin{align*}
3 \sqrt{\alpha}\,(x' + \rho^2 y' + \rho z') &= - \sqrt{3 \alpha}\,(1 + 2\rho)(x + \rho y + \rho^2 z),\\
3 \sqrt{\alpha}\,(x' + \rho y' + \rho^2 z') &= \sqrt{3 \beta}\,(1 + 2\rho)(x + \rho^2 y + \rho z).
\end{align*}

\medskip

The preceding investigations have been in my possession for about eighteen months.

\medskip
\noindent\textit{2 Stone Buildings, April 18, 1855.}

%==============================================================================
\section*{130.}
\addcontentsline{toc}{section}{130. Deuxieme memoire sur les fonctions doublement periodiques}
\subsection*{Deuxi\`eme M\'emoire sur les Fonctions Doublement P\'eriodiques.}

\selectlanguage{french}

[From the \textit{Journal de Math\'ematiques Pures et Appliqu\'ees} (Liouville), tom.\ xix.\ (1854), pp.\ 193--208: Sequel to Memoir t.~x.\ (1845), \textbf{25}.]

\medskip

Je vais essayer de d\'evelopper ici les propri\'et\'es qui se rapportent aux transformations lin\'eaires des p\'eriodes des fonctions $yx$, $gx$, $\Omega x$, $Zx$, dont je me suis occup\'e dans le M\'emoire sur les fonctions doublement p\'eriodiques que j'ai donn\'e dans ce Recueil en 1845. Avant d'entrer en mati\`ere, je remarque que, partant des expressions
\begin{align*}
\Omega &= \omega + \omega' i,\\
\mathrm{T} &= v + v' i,
\end{align*}
des deux p\'eriodes, o\`u $i = \sqrt{-1}$, on obtient, en \'ecrivant
\begin{align*}
\Omega^* &= \omega - \omega' i,\\
\mathrm{T}^* &= v - v' i,
\end{align*}
les \'equations
\begin{align*}
\Omega\, \mathrm{T}^* &= \omega v + \omega' v' - i (\omega v' - \omega' v),\\
\Omega^*\mathrm{T} &= \omega v + \omega' v' + i (\omega v' - \omega' v);
\end{align*}
au moyen desquelles et des valeurs
\[
B = \frac{\pi (\omega v' - \omega' v)}{\Omega\,\mathrm{T}\,\md(\omega v' - \omega' v)},\quad
\beta = \frac{\pi (\omega v + \omega' v')}{\Omega\,\mathrm{T}\,\md(\omega v' - \omega' v)}
\]
des quantit\'es $B$, $\beta$, on d\'eduit les formules
\[
B + \beta = \frac{\pi}{\Omega\,\md(\omega v' - \omega' v)}\,\Omega^*,\qquad
B - \beta = \frac{\pi}{\mathrm{T}\,\md(\omega v' - \omega' v)}\,\mathrm{T}^*.
\]

\medskip

Je ne fais attention qu'aux transformations qui correspondent \`a des entiers impairs et premiers, et je suppose, de plus, que la transformation soit toujours propre et r\'eguli\`ere, c'est-\`a-dire qu'en \'ecrivant
\[
(2k+1)\,\Omega_{,} = \lambda\Omega + \mu\mathrm{T} = (\lambda, \mu),
\]
\[
(2k+1)\,\mathrm{T}_{,} = \nu\Omega + \rho\mathrm{T} = (\nu, \rho),
\]
o\`u $2k+1$ est un entier positif, impair et premier, et o\`u $\lambda$, $\mu$, $\nu$, $\rho$ sont des entiers tels, qu'au signe pr\`es, $\lambda\rho - \mu\nu$ soit \'egal \`a $2k+1$, je suppose
\[ \lambda\rho - \mu\nu = 2k+1, \]
(condition pour que la transformation soit propre), et, en outre,
\begin{align*}
\lambda &\equiv 1,\quad \mu \equiv 0, \pmod{2}\\
\nu &\equiv 0,\quad \rho \equiv 1,
\end{align*}
(condition pour que la transformation soit r\'eguli\`ere).

\medskip

On trouve tout de suite
\[
\Omega = \rho\Omega_{,} - \mu\mathrm{T}_{,} = (\rho, -\mu),_{,}
\]
\[
\mathrm{T} = -\nu\Omega_{,} + \lambda\mathrm{T}_{,} = (-\nu, \lambda),_{,}
\]
j'\'ecris aussi
\[
\Omega_{,}^{*} = \omega_{,} - \omega'_{,}\,i,\quad \mathrm{T}_{,}^{*} = v_{,} - v'_{,}\,i,
\]
et je suppose que $B_{,}$, $\beta_{,}$ soient des fonctions de $\omega_{,}$, $v_{,}$ telles que les fonctions $B$, $\beta$ le sont de $\omega$, $v$.

\medskip

Cela \'etant, je forme d'abord l'\'equation
\[
(2k+1)(\omega_{,} v'_{,} - \omega'_{,} v_{,}) = \omega v' - \omega' v,
\]
au moyen de laquelle l'\'equation
\[
(B + \beta)\,\Omega = \frac{\pi}{\md(\omega v' - \omega' v)}\,\Omega^*
\]
se transforme en
\[
(2k+1)(B_{,} + \beta_{,})\,\Omega = \frac{\pi}{\md(\omega v' - \omega' v)}\,\Omega_{,}^{*}.
\]

\medskip

De l\`a
\[
\bigl[(2k+1)(B_{,} + \beta_{,}) - B\bigr]\,\Omega = \frac{\pi}{\Omega}\,\xi(\omega, \lambda),
\]
ou enfin
\[
\bigl[(2k+1)(B_{,} + \beta_{,}) - B\bigr]\,\Omega = -\frac{\pi}{\Omega}\,\beta_{,}(\omega, \mu),
\]
et de m\^eme
\[
\bigl[(2k+1)(B_{,} - \beta_{,}) - B_{,}\bigr]\,\Omega = -\frac{\pi}{\Omega}\,\xi_{,}(\omega, \lambda);
\]
\'equations qui seront bient\^ot utiles.

\medskip

Je suppose d'abord que $2k+1$ soit \'egal \`a l'unit\'e, transformation que l'on peut nommer triviale. La fonction $yx$ est d\'efinie par l'\'equation
\[
yx = e^{-\frac{1}{2}Bx^2} x \prod \Bigl[1 + \frac{x}{(m, n)}\Bigr],\quad \md(m, n) < T,\ T = \infty;
\]
dans $(m, n) = m\Omega + n\mathrm{T}$, les entiers $m$, $n$ doivent prendre toutes les valeurs positives ou n\'egatives (le seul syst\`eme $m = 0$, $n = 0$ except\'e) qui satisfont \`a l'in\'egalit\'e
\[ \md(m, n) < T, \]
dont le second membre $T$ sera ensuite suppos\'e infini. Soit $y_{,} x$ la fonction correspondante pour les p\'eriodes $\Omega_{,}$, $\mathrm{T}_{,}$, on aura
\[
y_{,} x = e^{-\frac{1}{2} B_{,} x^2} x \prod \Bigl[1 + \frac{x}{(m, n)_{,}}\Bigr],\quad \md(m, n)_{,} < T,\ T = \infty.
\]

\medskip

\noindent Or
\begin{align*}
(m, n)_{,} &= m\Omega_{,} + n\mathrm{T}_{,}\\
&= m(\lambda\Omega + \mu\mathrm{T}) + n(\nu\Omega + \rho\mathrm{T})\\
&= (\lambda m + \nu n)\,\Omega + (\mu m + \rho n)\,\mathrm{T}\\
&= m_{,}\Omega + n_{,}\mathrm{T}\\
&= (m_{,}, n_{,}).
\end{align*}

En \'ecrivant, comme nous venons de le faire,
\begin{align*}
m_{,} &= \lambda m + \nu n,\\
n_{,} &= \mu m + \rho n,
\end{align*}
on voit tout de suite qu'\`a chaque syst\`eme de valeurs enti\`eres de $m$, $n$, correspond un syst\`eme, et un seul syst\`eme, de valeurs enti\`eres de $m_{,}$, $n_{,}$; et que de chaque syst\`eme de valeurs enti\`eres de $m_{,}$, $n_{,}$, correspond un syst\`eme, et un seul syst\`eme, de valeurs enti\`eres de $m$, $n$; de plus, les syst\`emes $m = 0$, $n = 0$ et $m_{,} = 0$, $n_{,} = 0$, correspondent l'un \`a l'autre. Il est donc permis d'\'ecrire
\[
\prod \Bigl[1 + \frac{x}{(m, n)}\Bigr] = \prod \Bigl[1 + \frac{x}{(m, n)_{,}}\Bigr],
\]
les limites comme auparavant; car, \`a cause de
\[ (m, n)_{,} = (m_{,}, n_{,}), \]
la condition pour les limites, savoir
\[ \md(m, n)_{,} < T,\ T = \infty, \]
devient
\[ \md(m, n) < T,\ T = \infty. \]

\medskip

Cela donne enfin l'\'equation
\[
y_{,} x = e^{-\frac{1}{2}(B_{,} - B) x^2}\, yx,
\]
et, au moyen de cette \'equation, on obtient une \'equation correspondante pour la transformation de l'une quelconque des fonctions $yx$, $gx$, $\Omega x$, $Zx$, d\'efinies par les \'equations
\begin{align*}
yx &= e^{-\frac{1}{2}B x^2} x \prod \Bigl[1 + \frac{x}{(m, n)}\Bigr],\quad \md(m, n) < T,\\
gx &= e^{-\frac{1}{2}B x^2} \prod \Bigl[1 + \frac{x}{(m, n)}\Bigr],\quad \md(m, n) < T,\quad T = \infty,\\
\Omega x &= e^{-\frac{1}{2}B x^2} \prod \Bigl[1 + \frac{x}{(m, n)}\Bigr],\quad \md(m, n) < T,\\
Zx &= e^{-\frac{1}{2}B x^2} \prod \Bigl[1 + \frac{x}{(m, n)}\Bigr],\quad \md(m, n) < T,
\end{align*}
(\'equations dans lesquelles $m = m + \tfrac{1}{2}$, $n = n + \tfrac{1}{2}$). Je prends par exemple la fonction $gx$, et j'\'ecris dans l'\'equation entre $y_{,} x$ et $yx$, $x + \tfrac{1}{2}\Omega$, au lieu de $x$. Soit pour un moment $\rho = 2\rho' + 1$, $\mu = 2\mu'$; cela donne
\[
x + \tfrac{1}{2}\Omega = x + \tfrac{1}{2}(\rho\Omega_{,} - \mu\mathrm{T}_{,}) = x + (\rho', -\mu')_{,}.
\]

Donc
\[
y_{,}(x + \tfrac{1}{2}\Omega) = e^{\frac{1}{2}B_{,}(\rho', -\mu')_{,}^{2}}\, M_{,}\, g_{,} x,
\]
c'est-\`a-dire
\[
y_{,}(x + \tfrac{1}{2}\Omega_{,}) = e^{\frac{1}{2}B_{,}(\rho', -\mu')_{,}^{2}}\, M_{,}\, g_{,} x;
\]
de plus,
\[
y(x + \tfrac{1}{2}\Omega) = e^{\frac{1}{2}B \Omega^2}\, M\, g x.
\]

Ces substitutions \'etant effectu\'ees, les coefficients $M$, $M_{,}$ doivent \^etre \'elimin\'es en \'ecrivant $x = 0$; cela donne
\[
g_{,} x = e^{-\frac{1}{2}(B_{,} - B) x^2 \cdot \text{etc.}}\, g x,
\]
ou enfin, au moyen d'une \'equation d\'ej\`a trouv\'ee,
\[
g_{,} x = e^{-\frac{1}{2}(B_{,} - B) x^2}\, g x,
\]
et de m\^eme pour les fonctions $\Omega x$, $Zx$.

\medskip

Donc enfin, en repr\'esentant par $Jx$ l'une quelconque des fonctions $yx$, $gx$, $\Omega x$, $Zx$, on aura
\[
J_{,} x = e^{-\frac{1}{2}(B_{,} - B) x^2}\, Jx,
\]
o\`u $J_{,} x$ est ce que devient $Jx$ au moyen d'une transformation triviale (propre et r\'eguli\`ere) des p\'eriodes.

\medskip

Je passe \`a pr\'esent \`a la transformation pour un nombre impair et premier $(2k+1)$ quelconque; mais pour cela on a besoin de connaitre la valeur de la fonction
\[
u' = \prod \Bigl[1 + \frac{x}{(m, n) + y}\Bigr],\quad \md\{(m, n) + y\} < T,\ T = \infty,
\]
o\`u $y = a + bi$ est une quantit\'e r\'eelle ou imaginaire quelconque.

\medskip

Soit $u$ ce que devient $u'$ en prenant pour la condition par rapport aux limites
\[ \md(m, n) < T,\ T = \infty; \]
on trouve sans peine
\[
u = \frac{x + y}{y}\cdot\frac{y(x+y)}{y(y)}.
\]

\medskip

Pour trouver $u'$, je forme l'\'equation
\[
u : u' = \prod\Bigl[1 + \frac{x}{(m, n)}\Bigr],
\]
la limite inf\'erieure du produit infini double \'etant
\[ \md\{(m, n) + y\} > T, \]
et la limite sup\'erieure
\[ \md(m, n) < T,\ T = \infty; \]
cela donne
\[
\log u - \log u' = x \sum \frac{1}{(m, n)^2} - \tfrac{1}{2} x^2 \sum \frac{1}{(m, n)^3} + \cdots,
\]
car on peut d\'emontrer que
\[
\sum \frac{1}{(m, n)^k} = 0,\quad k \geq 2.
\]
Pour cela, observons que $m$ et $n$ \'etant infinis puisque $T$ l'est, la premi\`ere des sommes dont il s'agit peut se remplacer par l'int\'egrale double
\[
I = \iint \frac{dm\,dn}{(m, n)^2},
\]
laquelle (en \'ecrivant $m = r \cos\theta$, $n = r \sin\theta$, ce qui donne, comme on sait, $dm\,dn = r\,dr\,d\theta$) devient
\[
I = \iint \frac{dr\,d\theta}{r(\Omega \cos\theta + \mathrm{T}\sin\theta)^2},
\]
d'o\`u
\[
I = \iint \frac{(\log r)\,d\theta}{(\Omega \cos\theta + \mathrm{T}\sin\theta)^2},
\]
en prenant $(\log r)$ entre les limites convenables. Pour trouver ces limites, j'\'ecris
\[
(m, n) + y = r (\Omega \cos\theta + \mathrm{T}\sin\theta) + y,
\]
ce qui donne
\[
\md^2\{(m, n) + y\} = \bigl\{r(\Omega \cos\theta + \mathrm{T}\sin\theta) + y\bigr\}\bigl\{r(\Omega^* \cos\theta + \mathrm{T}^*\sin\theta) + y^*\bigr\},
\]
savoir, \`a l'une des limites
\begin{align*}
& r^2 (\Omega \cos\theta + \mathrm{T}\sin\theta)(\Omega^* \cos\theta + \mathrm{T}^*\sin\theta)\\
&\quad + r\bigl\{y^*(\Omega \cos\theta + \mathrm{T}\sin\theta) + y(\Omega^* \cos\theta + \mathrm{T}^*\sin\theta)\bigr\} + y y^* = 0;
\end{align*}
ou, en n\'egligeant les puissances n\'egatives de $T$,
\[
r = \frac{T}{\sqrt{(\Omega \cos\theta + \mathrm{T}\sin\theta)(\Omega^* \cos\theta + \mathrm{T}^*\sin\theta)}}
- \tfrac{1}{2}\Bigl[\frac{y}{\Omega \cos\theta + \mathrm{T}\sin\theta} + \frac{y^*}{\Omega^* \cos\theta + \mathrm{T}^*\sin\theta}\Bigr],
\]
et \`a l'autre limite,
\[
r = \frac{T}{\sqrt{(\Omega \cos\theta + \mathrm{T}\sin\theta)(\Omega^* \cos\theta + \mathrm{T}^*\sin\theta)}}.
\]
Or, en repr\'esentant ces deux \'equations par
\[ r = R - \phi,\quad r = R, \]
on trouve, pour la valeur de $(\log r)$ entre les deux limites,
\[
\log R - \log (R - \phi) = -\log\Bigl(1 - \frac{\phi}{R}\Bigr) = 0,
\]
\`a cause de la valeur infinie de $R$. Ainsi la somme cherch\'ee est nulle; et il est tout clair que les sommes suivantes $\sum \dfrac{1}{(m, n)^k}$, \&c., se r\'eduisent de m\^eme \`a z\'ero.

\medskip

Donc enfin,
\[
\log u - \log u' = x \sum \frac{1}{(m, n)^2}.
\]
Cela fait voir que
\[
u' = e^{-kx} u,
\]
le coefficient $k$ \'etant donn\'e au moyen de l'\'equation
\[
k = \sum \frac{1}{(m, n)^2},
\]
o\`u la somme est prise, comme auparavant, entre les limites
\[
\md\{(m, n) + y\} > T,\quad \md(m, n) < T,\ T = \infty.
\]
Mais il n'est pas permis d'\'ecrire
\[
k = \iint \frac{dm\,dn}{(m, n)^2}.
\]
En effet, cette int\'egrale n'est que le premier terme d'une suite dont il faudrait, pour obtenir un r\'esultat exact, prendre deux termes; le second terme de la suite serait une int\'egrale prise le long d'un contour, et il serait, ce me semble, tr\`es-difficile d'en trouver la valeur. Pour trouver la valeur de $k$, je remarque que $k$ sera fonction lin\'eaire des quantit\'es $T$, $y$, $y^*$, $\dfrac{1}{T}$, \&c., qui entrent dans les valeurs de $r$; donc, puisqu'en derni\`ere analyse $T = \infty$, $k$ ne peut \^etre que de la forme $L y + M y^*$. Cela \'etant, en substituant pour $u'$ sa valeur, je forme l'\'equation
\[
\frac{y(x + y)}{y(y)} = e^{-Bxy} e^{(-Bxy + Lxy + Mxy^*) x^2} \prod\Bigl[1 + \frac{x}{(m, n) + y}\Bigr]^{-1},\quad \md\{(m, n) + y\} < T,\ T = \infty,
\]
et j'\'ecris successivement
\[ y = \tfrac{1}{2}\Omega,\quad y = \tfrac{1}{2}\mathrm{T}, \]
ce qui donne pour les valeurs correspondantes du produit infini double $e^{-\frac{1}{2}Bx^2} g x$ et $e^{-\frac{1}{2}Bx^2} \Omega x$. En comparant les valeurs ainsi obtenues avec les \'equations qui donnent les valeurs de $y(x + \tfrac{1}{2}\Omega)$, $y(x + \tfrac{1}{2}\mathrm{T})$, on trouve
\[
L = 0,\quad M = \frac{\pi}{\md(\omega v' - \omega' v)};
\]
ou enfin,
\[
\frac{y(x + y)}{y(y)} = e^{-Bxy} e^{\beta x y^*} \prod\Bigl[1 + \frac{x}{(m, n) + y}\Bigr]^{-1},\quad \md\{(m, n) + y\} < T,\ T = \infty,
\]
laquelle est l'\'equation qu'il s'agissait d'\'etablir. Il est \`a peine n\'ecessaire de faire la remarque que pour $y = 0$, on doit consid\'erer \`a part le facteur $1 + \dfrac{x}{y}$, lequel multipli\'e par $y(y)$ devient tout simplement $x$; l'\'equation subsiste donc dans ce cas.

\medskip

En revenant au probl\`eme des transformations lin\'eaires, partant des \'equations
\begin{align*}
(2k+1)\,\Omega_{,} &= \lambda\Omega + \mu\mathrm{T},\\
(2k+1)\,\mathrm{T}_{,} &= \nu\Omega + \rho\mathrm{T},
\end{align*}
je suppose d'abord que les coefficients $\lambda$, $\nu$ ne satisfassent pas \`a la fois aux deux conditions
\[ \lambda \equiv 0,\quad \nu \equiv 0 \pmod{2k+1}, \]
et je prends $p$, $q$ des entiers quelconques tels, que $\lambda p + \nu q$ ne soit pas $\equiv 0$, $\md(2k+1)$.

\medskip

Cela \'etant, soient
\begin{align*}
\lambda p + \nu q &= p_{,},\\
\mu p + \rho q &= q_{,},\\
(2k+1)\,\psi &= p\Omega + q\mathrm{T},
\end{align*}
et, par cons\'equent,
\[ \psi = p\Omega_{,} + q\mathrm{T}_{,}. \]
Je forme l'\'equation
\[
(m_{,}, n_{,}) + s\psi = (m, n)_{,},
\]
savoir
\[
m_{,}\Omega + n_{,}\mathrm{T} + s\psi = m\Omega_{,} + n\mathrm{T}_{,},
\]
c'est-\`a-dire
\begin{align*}
\lambda m + \nu n - s p_{,} &= (2k+1)\,m_{,},\\
\mu m + \nu n - s q_{,} &= (2k+1)\,n_{,},
\end{align*}
ou, ce qui est la m\^eme chose,
\begin{align*}
m - sp &= m_{,}\rho - n_{,}\nu,\\
n - sq &= -m_{,}\mu + n_{,}\lambda.
\end{align*}
Or, $m_{,}$, $n_{,}$, $s$ \'etant des entiers donn\'es, $m$, $n$ seront aussi des entiers; de m\^eme, $m$, $n$ \'etant des entiers donn\'es, on trouve de $-k$ \`a $-k$ un entier $s$ qui donne $m_{,}$ un entier. Mais cela \'etant, $n_{,}$ sera aussi un entier; car autrement $n_{,}$ serait une fraction ayant pour d\'enominateur, lequel on voudrait, des nombres $2k+1$, $\lambda$, $\nu$, ce qui est impossible \`a moins que
\[ \lambda \equiv 0,\quad \nu \equiv 0 \pmod{2k+1}. \]
Mais si ces \'equations avaient lieu, on trouverait d'abord $s$ de mani\`ere \`a avoir $m_{,}$ entier, et alors, puisqu'on n'a pas aussi
\[ \mu \equiv 0,\quad \rho \equiv 0 \pmod{2k+1} \]
(en effet, cela est impossible \`a cause de l'\'equation $\lambda\rho - \mu\nu = 2k+1$), on d\'emontrerait, comme auparavant, pour $n_{,}$, que $m_{,}$ est entier. Donc, enfin, $m$, $n$ \'etant des entiers donn\'es, on trouve pour $m_{,}$, $n_{,}$, $s$ un syst\`eme d'entiers tel que $s$ soit compris de $k$ \`a $-k$, et l'on voit sans peine qu'il n'y a qu'un seul syst\`eme de cette esp\`ece.

\medskip

A pr\'esent, partant de l'\'equation
\[
\frac{y(x + y)}{y(y)} = e^{-Bxy + \beta x y^*}\prod\Bigl[1 + \frac{x}{(m, n) + y}\Bigr]^{-1},
\]
(et faisant attention \`a la particularit\'e que pr\'esente le cas de $y = 0$), j'\'ecris successivement
\[ y = 0,\quad y = \pm\psi,\ \ldots,\ y = \pm k\psi, \]
et je forme le produit des \'equations ainsi trouv\'ees. Cela donne, \`a cause de $(m_{,}, n_{,}) + s\psi = (m, n)_{,}$,
\[
\prod \frac{y(x + s\psi)}{y(s\psi)} = e^{-Bxy\,\text{etc.}} \prod\Bigl[1 + \frac{x}{(m_{,}, n_{,}) + s\psi}\Bigr]^{-1},
\]
la condition, par rapport aux limites, \'etant
\[ \md(m, n)_{,} < T,\ T = \infty. \]
Or
\[
y_{,} x = e^{-\frac{1}{2} B_{,} x^2} x \prod \Bigl[1 + \frac{x}{(m_{,}, n_{,})}\Bigr],
\]
avec la m\^eme condition, par rapport aux limites; donc, enfin,
\[
y_{,} x = \frac{e^{-\frac{1}{2} B_{,} x^2 + \cdots}}{e^{\cdots}}\prod \frac{y(x + s\psi)}{y(s\psi)},
\]
o\`u, dans le num\'erateur, $s$ doit avoir toutes les valeurs enti\`eres depuis $s = -k$ jusqu'\`a $s = +k$, $y$ compris $s = 0$, et dans le d\'enominateur ces m\^emes valeurs, hormis la valeur $s = 0$.

\medskip

Il est, \`a pr\'esent, facile de faire voir que cette propri\'et\'e subsiste pour l'une quelconque des fonctions $yx$, $gx$, $\Omega x$, $Zx$; en effet, pour la d\'emontrer pour $gx$, j'\'ecris $x + \tfrac{1}{2}\Omega$ au lieu de $x$; en prenant, pour un moment, $\rho = 2\rho' + 1$, $\mu = 2\mu'$; cela donne
\[
x + \tfrac{1}{2}\Omega = x + (\rho', -\mu')_{,}.
\]
Donc
\[
y_{,}(x + \tfrac{1}{2}\Omega) = e^{\frac{1}{2} B_{,}(\rho', -\mu')_{,}^{2}}\, M\, g_{,} x.
\]

\medskip

A pr\'esent, partant de l'\'equation
\[
\frac{y(x + y)}{y(y)} = e^{-Bxy + \beta x y^*}\prod\Bigl[1 + \frac{x}{(m, n) + y}\Bigr]^{-1}
\]
(en faisant attention \`a la particularit\'e que pr\'esente le cas de $y = 0$), j'\'ecris successivement
\[ y = 0,\quad y = \pm \psi,\ \ldots,\ y = \pm k\psi, \]
et je forme le produit des \'equations ainsi trouv\'ees. Cela donne, \`a cause de $(m_{,}, n_{,}) + s\psi = (m, n)_{,}$,
\[
\prod \frac{y(x + s\psi)}{y(s\psi)} = e^{-Bxy + \cdots} \prod\Bigl[1 + \frac{x}{(m_{,}, n_{,}) + s\psi}\Bigr]^{-1},
\]
la condition, par rapport aux limites, \'etant
\[ \md(m, n)_{,} < T,\ T = \infty. \]

\medskip

\noindent (\emph{D'autres manipulations alg\'ebriques analogues conduisent aux formules concr\`etes pour la transformation des fonctions $yx$, $gx$, $\Omega x$, $Zx$, exprim\'ees par produits sur les indices $s = -k, \ldots, k$ et $s = -k, \ldots, k$ avec $s \neq 0$ au d\'enominateur. Les formules d\'efinitives s'\'ecrivent comme suit:})

\medskip

\[
\Omega_{,} = (\Omega - \psi)\,\mathrm{T}_{,} = (\mathrm{T} - \psi),\ \text{etc.,}
\]
et, en r\'esumant, pour les fonctions transform\'ees, on obtient des syst\`emes du type
\begin{align*}
\Omega_{,} &= \Omega_{,1}(2\Omega + 2\mathrm{T}_{,}),\\
\mathrm{T}_{,} &= \Omega_{,1}\Omega + \mathrm{T}_{,1}\mathrm{T} = \psi,\\
\mathrm{T}_{,2} &= \mathrm{T}.
\end{align*}
Les param\`etres particuliers correspondent au cas g\'en\'eral en fonctions de $\Omega$, $\mathrm{T}$; on obtient \'egalement, par un moyen d'une transformation triviale, on obtiendrait
\[ B' = \Omega + 2 B_{,}\, \mathrm{T},\quad \mathrm{T}_{,} = \mathrm{T}, \]
et puis
\[
\psi = \frac{1}{2k+1}\,\Omega,\quad \mathrm{T}_{,} = \mathrm{T},
\]
et, de l\`a,
\[
\frac{\Omega_{,}}{\mathrm{T}_{,}} = \frac{1}{2k+1}\cdot\frac{\Omega}{\mathrm{T}}.
\]
Les \'equations correspondantes pour la deuxi\`eme sont:
\begin{align*}
p &= 2k+1,\quad q = 0,\\
\lambda &= 2k+1,\quad \nu = 0,\\
\mu &= 0,\quad \rho = 1,\\
\rho_{,} &= 1,\quad q_{,} = 0,\\
\psi_{,} &= \frac{1}{2k+1}\bigl[(2k+1)\,\Omega + (2k+1)\,\mathrm{T}\bigr],\\
\Omega_{,1} &= \Omega,\\
\mathrm{T}_{,1} &= (2k+1)\,\Omega + \mathrm{T},
\end{align*}
ce qui donne
\[
\frac{\Omega_{,}}{\mathrm{T}_{,}} = (2k+1)\,\frac{\Omega}{\mathrm{T}}.
\]

\medskip

J'ajoute, sans m'arr\^eter pour les d\'emontrer, quelques formules de transformation pour le nombre $3$; je trouve d'abord
\begin{align*}
\Omega_{,} &= \tfrac{1}{3}\Omega,\quad \mathrm{T}_{,} = \mathrm{T},\\
B_{,} &= e^{-\frac{1}{2}B\Omega^2 - \beta\Omega\Omega^*}\,\frac{y(x + \tfrac{1}{3}\Omega)\,y(x - \tfrac{1}{3}\Omega)}{y^2(\tfrac{1}{3}\Omega)},\\
G_{,} &= e^{-\frac{1}{2}B\Omega^2 - \beta\Omega\Omega^*}\,\frac{y(x + \tfrac{1}{3}\Omega)\,y(x - \tfrac{1}{3}\Omega)}{B^2(\tfrac{1}{3}\Omega)},\\
Z_{,} &= e^{-\frac{1}{2}B\Omega^2 - \beta\Omega\Omega^*}\,\frac{y(x + \tfrac{1}{3}\Omega)\,y(x - \tfrac{1}{3}\Omega)}{B^2(\tfrac{1}{3}\Omega)}.
\end{align*}

Ces \'equations donnent, en introduisant les fonctions elliptiques $\phi x$, $fx$, $Fx$ donn\'ees au moyen de
\[
\phi x = \frac{Yx}{Bx},\quad fx = \frac{gx}{Bx},\quad Fx = \frac{Gx}{Bx},
\]
les \'equations
\begin{align*}
\phi_{,} x &= \phi x\cdot\frac{f x\cdot f(x + \tfrac{1}{3}\Omega)\, f(x - \tfrac{1}{3}\Omega)}{f^2(\tfrac{1}{3}\Omega)},
\end{align*}
dont la seconde peut s'\'ecrire sous la forme
\[
f_{,} x = \frac{1 - x^2\,\dfrac{P}{Q}(\tfrac{1}{3}\Omega)\,\phi x}{1 + x^2 \dfrac{P}{Q}(\tfrac{1}{3}\Omega)\,\phi x},\quad \text{etc.}
\]
On trouve en effet, en mettant comme \`a l'ordinaire $\Phi = \phi^2 x$,
\begin{align*}
c'^2 &= b\phi c'\,,\quad
e'^2 = (b - c')\,c'\,,
\end{align*}
et puis
\begin{align*}
\phi_{,} x &= \frac{1 - c(1 - b)\,\phi x}{1 - c(c - b)\,\phi x},\\
f_{,} x &= \frac{1 - (c + b)\,\phi x}{1 - c(c - b)\,\phi x},\\
F_{,} x &= \frac{1 - c\phi x}{1 - c(c - b)\,\phi x},
\end{align*}
formules qui correspondent \`a celles de la transformation de Lagrange. Les \'equations $(\beta)$, $(\gamma)$ donnent en mettant ces valeurs de $\phi$, $f$, $F$, laquelle, \'egal\'ee \`a la valeur qui vient d'\^etre trouv\'ee, donne
\[
\frac{xc\,(\overline{\phi}\,P(\tfrac{1}{3}\Omega))}{B(x - \tfrac{1}{3}\Omega)\,B(x + \tfrac{1}{3}\Omega)} = \frac{\phi x\,fx}{1 - c(c - b)\,\phi x}.
\]

\medskip

\noindent (\textit{Suit la d\'erivation par formules analogues du syst\`eme correspondant pour la transformation $\Omega_{,} = \Omega$, $\mathrm{T}_{,} = \tfrac{1}{3}\mathrm{T}$. Le texte expose ensuite le syst\`eme correspondant, puis poursuit avec les formules d\'efinitives.})

\medskip

On obtient tout de suite, pour le premier cas, le syst\`eme d'\'equations
\begin{align*}
\rho_{,} &= 1,\quad q_{,} = 2q_{,},\\
\lambda &= 1,\quad \nu = 2,\\
\mu &= 0,\quad \rho = (2k+1),\\
p &= 1,\quad q = 1,\\
\psi &= \tfrac{1}{2k+1}\bigl[(2k+1) + 2\,\mathrm{T}_{,}\bigr]\,\mathrm{T},\\
\Omega_{,} &= \Omega,\\
\mathrm{T}_{,} &= (2k+1)\,\mathrm{T} + \psi,\\
\mathrm{T}_{,2} &= \mathrm{T}.
\end{align*}

Le tout particulier le plus simple est celui de $q = 0$, cela donne
\[ \psi = \Omega_{,},\quad \mathrm{T}_{,} = \mathrm{T}; \]
et, de l\`a,
\[
\frac{\Omega_{,}}{\mathrm{T}_{,}} = \frac{1}{2k+1}\cdot\frac{\Omega}{\mathrm{T}},
\]
et m\^eme le cas g\'en\'eral se r\'eduit \`a celui-ci, car, au moyen d'une transformation triviale, on obtiendrait
\[
B' = \Omega + 2 B_{,}\,\mathrm{T},\quad \mathrm{T}_{,} = \mathrm{T},
\]
et puis
\[
\psi = \frac{1}{2k+1}\Omega',\quad \mathrm{T}_{,} = \mathrm{T},
\]
et, de l\`a,
\[
\frac{\Omega_{,}}{\mathrm{T}_{,}} = \frac{1}{2k+1}\cdot\frac{\Omega}{\mathrm{T}}.
\]

\medskip

Les \'equations correspondantes pour la deuxi\`eme sont:
\begin{align*}
p &= 2k+1,\quad q = 0,\\
\lambda &= 2k+1,\quad \nu = 0,\\
\mu &= 0,\quad \rho = 1,\\
p_{,} &= 1,\quad q_{,} = 0,\\
\psi &= \tfrac{1}{2k+1}\bigl[(2k+1)\,\Omega + (2k+1)\,\mathrm{T}\bigr],\\
\Omega_{,1} &= \Omega,\\
\mathrm{T}_{,1} &= (2k+1)\,\Omega + \mathrm{T},
\end{align*}
ce qui donne
\[
\frac{\Omega_{,}}{\mathrm{T}_{,}} = (2k+1)\,\frac{\Omega}{\mathrm{T}}.
\]

\medskip

J'ajoute, sans m'arr\^eter pour les d\'emontrer, quelques formules de transformation pour le nombre $3$; je trouve d'abord
\begin{align*}
\Omega_{,} &= \tfrac{1}{3}\Omega,\quad \mathrm{T}_{,} = \mathrm{T},\\
y_{,} x &= e^{-\frac{1}{2}\cdots x^2}\, gx\,Ex,\\
g_{,} x &= e^{-\frac{1}{2}\cdots x^2}\,\frac{E(x + \tfrac{1}{3}\Omega)\,E(x - \tfrac{1}{3}\Omega)}{E^2(\tfrac{1}{3}\Omega)},\\
\Omega_{,} x &= e^{-\frac{1}{2}\cdots x^2}\,\frac{E(x + \tfrac{1}{3}\Omega - \tfrac{1}{2}\mathrm{T})\,E(x - \tfrac{1}{3}\Omega - \tfrac{1}{2}\mathrm{T})}{E^2(\tfrac{1}{3}\Omega - \tfrac{1}{2}\mathrm{T})},\\
Z_{,} x &= e^{-\frac{1}{2}\cdots x^2}\,\frac{E(x + \tfrac{1}{3}\Omega - \tfrac{1}{2}\mathrm{T})\,E(x - \tfrac{1}{3}\Omega - \tfrac{1}{2}\mathrm{T})}{E^2(\tfrac{1}{3}\Omega - \tfrac{1}{2}\mathrm{T})}.
\end{align*}

Ces \'equations donnent, en introduisant les fonctions elliptiques $\phi x$, $fx$, $Fx$ donn\'ees ci-dessus, les \'equations
\begin{align*}
\phi_{,} x &= \phi x\cdot\frac{f(x + \tfrac{1}{3}\Omega)\,f(x - \tfrac{1}{3}\Omega)}{f^2(\tfrac{1}{3}\Omega)},\\
f_{,} x &= \frac{1 - c\phi x}{1 - c(c - b)\,\phi x},\\
F_{,} x &= \frac{1 - c\phi x}{1 - c(c - b)\,\phi x};
\end{align*}
\'equations qui sont, sauf la notation dont je me suis d\'ej\`a servi dans la m\'emoire cit\'e, les m\^emes que celles de Jacobi.

\medskip
\noindent\textit{Londres, Stone Buildings, 23 F\'ev.\ 1852.}

\selectlanguage{english}

%==============================================================================
\section*{131.}
\addcontentsline{toc}{section}{131. Nouvelles recherches sur les covariants}
\subsection*{Nouvelles Recherches sur les Covariants.}

\selectlanguage{french}

[From the \textit{Journal f\"ur die reine und angewandte Mathematik} (Crelle), tom.\ xlvii.\ (1854), pp.\ 109--125.]

\medskip

Je me sers de la notation
\[
(a_0, \ldots a_n \mid x, y)^n
\]
pour repr\'esenter la fonction
\[
a_0 x^n + \tfrac{n}{1}\, a_1 x^{n-1} y + \cdots + a_n y^n;
\]
en supposant que les coefficients $a'_0, a'_1, \ldots a'_n$ soient donn\'es par l'\'equation
\[
(a'_0, \ldots a'_n \mid \lambda x + \mu y, \lambda' x + \mu' y)^n = (a_0, a_1, \ldots a_n \mid x, y)^n,
\]
suppos\'e identique par rapport \`a $x$, $y$, soit $\phi(a_0, a_1, \ldots a_n; x, y)$ une fonction des coefficients et des variables, telle que
\[
\phi(a'_0, a'_1, \ldots a'_n; x, y) = (\lambda\mu' - \lambda'\mu)^a\,\phi(a_0, a_1, \ldots a_n; x, y),
\]
cette fonction $\phi$ est nomm\'ee dans M.\ Cauvarian, et dans le cas particulier de $\phi$ une fonction des seuls coefficients, un \emph{Invariant} de la fonction donn\'ee.

\medskip

Je suppose d'abord que les nouveaux coefficients soient donn\'es par l'\'equation
\[
(a'_0, \ldots a'_n \mid x, y)^n = \phi(a_0, a_1, \ldots a_n; x, y);
\]
cela suppose les relations
\begin{align*}
a'_0 &= a_0,\\
a'_1 &= a_1 + \lambda a_0,\\
a'_2 &= a_2 + 2\lambda a_1 + \lambda^2 a_0,\\
&\;\,\&\mathrm{c}.
\end{align*}

\medskip

Il faut donc que le covariant $\phi$ satisfasse \`a l'\'equation
\[
\phi(a'_0, a'_1, \ldots a'_n; x, y) = \phi(a_0, a_1, \ldots a_n; x, y + \lambda x),
\]
ce qui donne
\[
(a_0, a_1, \ldots a_n; x + py + p^2 z) = \phi(a_0, a_1, \ldots a_n; x, y),
\]
o\`u en posant
\begin{align*}
a''_0 &= a_0,\\
a''_1 &= a_1 + \lambda a_0,\\
a''_2 &= a_2 + 2\lambda a_1 + \lambda^2 a_0,\\
&\;\,\&\mathrm{c.},
\end{align*}
le covariant $\phi$ doit satisfaire aux \'equations
\[
\phi(a''_0, a''_1, \ldots a''_n; x, y) = \phi(a_0, a_1, \ldots a_n; x, y).
\]

\medskip

\noindent (\emph{La suite du m\'emoire d\'eveloppe la th\'eorie des covariants en partant de cette d\'efinition: \'etude des op\'erateurs $\square$, $\nabla$ et de leurs propri\'et\'es alg\'ebriques, \'enonc\'e du th\'eor\`eme fondamental sur la repr\'esentation des covariants par produits de racines, et applications aux invariants des formes binaires de degr\'es bas. Le texte n'\'etant constitu\'e que de formules denses et d'identit\'es alg\'ebriques avec de nombreux indices, nous reproduisons ci-dessous les principales \'equations et \'enonc\'es retenus du texte original.})

\medskip

Et r\'eciproquement, toute fonction homog\`ene par rapport aux coefficients et issue par rapport aux variables, qui satisfait \`a ces \'equations, est un covariant de la fonction donn\'ee.

\medskip

Par exemple, l'invariant $\phi = ac - b^2$ de la fonction $ax^2 + 2bxy + cy^2$ satisfait aux \'equations
\[
(ab + 2bk)\,k - (bk + ck)^2 = (2bk - a^2k)\,k - 0,
\]
et l'invariant $\phi = ace - bd^2 + abe^2 - 3 bcd - b^3 + 2ad^3$ de la fonction $ax^4 + 4 b x^3 y + 6 c x^2 y^2 + 4 d x y^3 + e y^4$ satisfait aux \'equations
\[
(a k_{0} + 3 b k)\, k_{1} = (a k_{0}, k_{1})\,k_{0}.
\]

\medskip

Il est clair qu'on ne consid\'erait que les fonctions qui restent les m\^emes en \'echange dans les deux indices la coefficient $a_0$, $a_1, \ldots a_n$ et les variables $x$, $y$, respectivement. Le nombre des coefficients $a_0$, $a_1, \ldots a_n$ est sup\'erieur, le nombre des coefficients ind\'etermin\'es; il est exactement \'egal au nombre des coefficients dans la fonction donn\'ee. Mais il y a quelque autre indication \`a faire avant de pouvoir consid\'erer comme arbitraires les indices $\beta$, $\&$c., dont la s\'erie indique le caract\`ere du covariant cherch\'e.

\medskip

Cherchons par exemple pour la fonction $ax^4 + 4 b x^3 y + 6 c x^2 y^2 + 4 d x y^3 + e y^4$ un covariant $\phi$ de la forme
\[
\phi = A x^4 y + B b x^3 y^2 + C x^2 y^3 + D y^4,
\]
contenant les quatre coefficients ind\'etermin\'es $A$, $B$, $C$, $D$. En substituant dans l'\'equation les m\^emes apr\`es substitutions convenables, on obtient
\[
3 (a A) x^2 d y^4 + (4 b B + 4 c C) x^2 y^3 + (4 c B + 4 d D) x y^3 + (8 e) y^4 = 0,
\]
ou bien, puisque cette \'equation doit avoir lieu ind\'ependamment de la quantit\'e $k$ (qui est arbitraire), il faut que les coefficients soient s\'epar\'ement nuls. On a donc les quatre \'equations
\[
a A = 0,\quad 4 b B = 0,\quad 4 c B = 0,\quad 8 e = 0,
\]
ce qui donne le syst\`eme $A : B : C : D$, et l'on voit que cette s\'erie est compl\`etement d\'etermin\'ee, et l'on obtient en effet
\[
\phi = a c + b d - 6 a c d - 4 e b^2 + 3 b c^2.
\]

\medskip

La circonstance suivante r\'eclame n\'eanmoins de l'attention. Il s'agit d'\'enoncer que dans le calcul, le probl\`eme de trouver le nombre des coefficients d'un cota donn\'e; probl\`eme qui n'a aucun lieu de l'autre.

\medskip

\noindent (\emph{La suite d\'eveloppe les notations symboliques $\Omega$, $\square$, etc., et culmine avec l'\'enonc\'e du th\'eor\`eme central de l'article:})

\medskip

\noindent\textbf{Th\'eor\`eme.}\quad Tout covariant $\phi$ de la fonction
\[ (a_0, \ldots a_n \mid x, y)^n \]
satisfait aux deux \'equations \`a diff\'erences partielles
\[
(\square - p)\,\phi = 0,\quad (\bigtriangleup - q)\,\phi = 0,
\]
o\`u
\begin{align*}
\square &= n a_1 \partial_{a_0} + (n-1)\,a_2 \partial_{a_1} + \cdots + a_n \partial_{a_{n-1}} + y \partial_x,\\
\bigtriangleup &= n a_0 \partial_{a_1} + (n-1)\,a_1 \partial_{a_2} + \cdots + a_{n-1}\partial_{a_n} + x \partial_y,
\end{align*}
et r\'eciproquement.

\medskip

Inversement on a, en notant d'abord $a^n = b$, puis en consid\'erant des trois quantit\'es ind\'ependantes $a$, $b$, $c$, la solution g\'en\'erale de l'\'equation \`a diff\'erences partielles
\[
ap\,\partial_a + 2bp\,\partial_b + 3cp\,\partial_c = 0,
\]
en supposant un poids de $p$, est ind\'ependante de $a$, $b$, $c$, et l'on consid\`ere ainsi cette propri\'et\'e fondamentale: \emph{tout covariant est compos\'e de termes de la premi\`ere esp\`ece}; et r\'esultat ainsi obtenu g\'en\'eral des covariants forment, lorsque cela est permis, par cette m\^eme observation, une suite indiquant la longueur des polyn\^omes \`a poids $p$.

\medskip

Pour les polyn\^omes binaires, l'\'enonc\'e du th\'eor\`eme prend la forme suivante:

\medskip

\noindent\textbf{Th\'eor\`eme.}\quad Pour qu'une fonction $\phi$ des coefficients $a_0, a_1, \ldots a_n$ d'une forme binaire de degr\'e $n$, soit un \emph{invariant}, il faut et il suffit qu'elle satisfasse simultan\'ement aux deux \'equations
\[
\square\,\phi = 0,\quad \bigtriangleup\,\phi = 0,
\]
o\`u $\square$, $\bigtriangleup$ sont les op\'erateurs ci-dessus, restreints \`a leur partie agissant sur les seuls coefficients.

\medskip

\noindent (\emph{Le texte poursuit avec des d\'emonstrations d\'etaill\'ees, des applications aux discriminants des formes du degr\'e $2$, $3$, $4$, et des exemples num\'eriques.})

\medskip

Pour le discriminant
\[
\nabla = ax^2 + 2bxy + cy^2
\]
de la fonction de degr\'e $2$, le discriminant est $4ac - b^2$. Nous savons alors
\[
\nabla = -b^2 (4ac - b^2) + B^2 + Cb^2,
\]
en faisant $F = 4ab + 2 b^3$, $B$, $C$ \'etant donn\'es par
\[
\eta\, B = F(4xy^2 - by^3),\quad 2yC = F(B),
\]
c'est-\`a-dire $B = 12 a c b - b^3$, $C = -27 a^2 b$, et de l\`a
\[
\nabla = -27 a^2 b + 18 a c b^2 - 4ac^3 + b^2 c^2 - 4 a^3 c^2,
\]
valeur qui correspond en effet au discriminant ordinaire. En changeant d'une mani\`ere convenable les coefficients.

\medskip
\noindent\textit{Londres, Stone Buildings, 25 F\'evr.\ 1852.}

\selectlanguage{english}

%==============================================================================
\section*{132.}
\addcontentsline{toc}{section}{132. Reponse a une question proposee par M. Steiner}
\subsection*{R\'eponse \`a une Question Propos\'ee par M. Steiner.}

\selectlanguage{french}

(Aufgabe 4, \textit{Crelle} t.\ xxxi.\ (1846) p.\ 90.)

\medskip

[From the \textit{Journal f\"ur die reine und angewandte Mathematik} (Crelle), tom.\ x.\ (1855), pp.\ 277--278.]

\medskip

En partant des deux th\'eor\`emes:

\medskip

I.\quad Qu'il existe un nombre infini de surfaces du second ordre qui touchent une surface du second ordre aux six ar\^etes d'un quelconque hexa\`edre quadrilat\'eral;

\medskip

II.\quad Que le lieu d'intersection de trois plans rectangles qui touchent une surface du second ordre est une sph\`ere concentrique avec la surface, tandis que pour les hyperbolo\"ides un seul nappe le lieu d'intersection est r\'eduit \`a un plan,

\medskip

M.\ Steiner suppose le cas d'un parall\'el\'epip\`ede rectangle, en raisonnant d'une c\^ote $P$ de ce parall\'el\'epip\`ede $P$; par lequel passent trois plans principaux. Les trois plans du parall\'el\'epip\`ede $P$ rencontrent deux surfaces ses faces $A$, $A'$, $B$, $B'$, $C$, $C'$. Avec ces six plans principaux $P$, $\&$c.\ on aurait six points de contact, qui pourraient former une figure du second ordre. Mais c'est dans ce cas, on doit avoir un point de l'espace consid\'er\'ee, $P$, dont les trois plans principaux \`a une surface dont les plans tangents $A$, $A'$, $B$, $B'$, $C$, $C'$ contiennent les six points et faces du parall\'el\'epip\`ede, et celles trois faces, sph\'eriques, et en outre, lorsque le parall\'el\'epip\`ede $P$ est aux trois plans principaux d'un point d'une surface dont les trois faces sont parall\`eles aux faces de $A'$, $B'$, $C'$. La contradiction qui apparaitra dans le cas g\'en\'eral est par cons\'equent admissible que pour les six points.

\medskip

Aussi d'attention le solution des \'equations $(A)$, je vais d\'emontrer quelque propri\'et\'e g\'en\'erale qui appartient \`a ce th\'eor\`eme. Pour abr\'eger, je consid\'ere $\Delta$ le syst\`eme principal de la surface; \\
ce dernier que les coefficients $a_0$, $a_1, \ldots a_n$ et les coefficients de la m\^eme forme donnent en effet le m\^eme nom et m\'eritera m\'eriter que les nouvelles indices ne sont pas $r$, mais $s$. Donc, tout autre que sont $a$, $b$, $c$, $d$, $\ldots$; on consid\`ere ensuite quelques points jusqu'au syst\`eme principal $P$, et l'on cherche les valeurs qu'elles donnent dans la fonction $a$, $b$, $c$, $d$, $\ldots$, $f$, qui en a l'expression. On consid\`ere alors les \'equations sym\'etriques de la fonction $a$ et de la fonction $f$, et l'on consid\`ere la quantit\'e $(\Phi - P^2 = 0)$, en effet, ce th\'eor\`eme est bien d\'etermin\'e, et le th\'eor\`eme $\beta$ et le II sont tous deux l\'egitimes.

\selectlanguage{english}

%==============================================================================
\section*{133.}
\addcontentsline{toc}{section}{133. Sur un theoreme de M. Schlafli}
\subsection*{Sur un Th\'eor\`eme de M. Schl\"afli.}

\selectlanguage{french}

[From the \textit{Journal f\"ur die reine und angewandte Mathematik} (Crelle), tom.\ l.\ (1855), pp.\ 278--282.]

\medskip

On lit dans le \S\ 13 d'un m\'emoire tr\`es int\'eressant de M.\ Schl\"affli intitul\'e ``\"Uber die Resultanten eines Systems mehrerer algebraischer Gleichungen,'' (\textit{M\'em.\ de l'Acad.\ de Vienne}, t.\ iv. [1852]) un tr\`es beau th\'eor\`eme sur les Z\'esultantes.

\medskip

Pour faire voir plus clairement, en quoi consiste ce th\'eor\`eme, je prends un cas particulier. Soit
\begin{align*}
U &= a x^2 + 3 b x^2 y + 3 c x y^2 + d y^3 + (a, b, c, d \mid x, y)^3,\\
V &= a x^3 + 2 b x y + c y^2 = (a, b, c \mid x, y)^2.
\end{align*}
Je fais $p = a^3$, $q = ay$, $r = pq$, et je forme les op\'erateurs
\begin{align*}
\nabla &= \tfrac{1}{2}\partial_a + p\partial_b + q\partial_c,\\
\bigtriangleup &= \tfrac{1}{2}\partial_a + p\partial_b + q\partial_c,
\end{align*}
lesquels, op\'erant sur $U$, donnent
\[
r (p\xi + qy - C),\quad p(p\xi + qy + C),
\]
op\'erant sur $V$, donnent
\[
p\xi + qy + r^2.
\]

\medskip

Cela \'etant, soit $\phi = 0$ le r\'esultant des \'equations $U = 0$, $V = 0$, c'est-\`a-dire l'\'equation que l'on obtient en \'eliminant $x$, $y$ entre les deux \'equations $U = 0$, $V = 0$, on a
\[
\phi = \det\begin{pmatrix}
a & 3b & 3c & d & \cdot\\
\cdot & a & 3b & 3c & d\\
a & 2b & c & \cdot & \cdot\\
\cdot & a & 2b & c & \cdot\\
\cdot & \cdot & a & 2b & c
\end{pmatrix}.
\]

Je suppose que les op\'erations $\nabla$, $\bigtriangleup$ op\`erent sur le r\'esultant $\phi$, et que cela donne les fonctions $\nabla\phi$, $\bigtriangleup\phi$, on trouve sans peine
\begin{align*}
\bigtriangleup\phi &= (p\xi + q\eta - C)\phi,\\
\nabla\phi &= (p\xi + q\eta + C)\phi,
\end{align*}
ou en \'crivant $(\xi, \eta, \zeta)$ pour les r\'esultats de $\phi$ et $(\xi)$ pour les m\^emes r\'esultats de $\phi$, on aurait \`a consid\'erer ce r\'esultant comme une fonction sym\'etrique de la fonction; ce r\'esultant est ce qui est nomm\'e dans le th\'eor\`eme de M.\ Schl\"affli.

\medskip

G\'en\'eralement, on suppose que l'on ait autant de fonctions $U$, $V$, $W$, $\ldots$ que d'ind\'etermin\'es $x$, $y$, $z$, $\ldots$, ces fonctions \'etant des degr\'es respectifs $\mu$, $\mu'$, \ldots, des nombres ind\'etermin\'es, et l'on pose pour chacune des fonctions du syst\`eme l'op\'erateur $\nabla$ corresp\^ond ce que ce que devient l'op\'erateur $\nabla$ pour la fonction $U$, dans laquelle on substitue \`a la place des ind\'etermin\'ees $x$, $y$, $z$, $W$, \ldots, des facteurs $p$, $q$, $r$, \&c.; avec les indices respectifs des fonctions $U$, $V$, $W$, \ldots, et de la fonction d\`ene faite, on a un th\'eor\`eme dont je ne pourrais indiquer la d\'emonstration de M.\ Schl\"affli, mais dont la valeur du r\'esultant $\phi$ est telle que cela est en outre une fonction sym\'etrique.

\medskip

Cela \'etant, soit $\phi$ le r\'esultant des fonctions $U$, $V$, $W$, $\ldots$, c'est-\`a-dire l'\'equation que l'on obtient en \'eliminant $x$, $y$, $z$, \ldots\ ; et soit $r$ le th\'eor\`eme de l'op\'eration $\nabla$ comme l'on a vu pour les hyperbolo\"ides. Cela donne en effet pour les r\'esultants \`a fonction de $\nabla$.

\medskip

\noindent\textbf{Th\'eor\`eme.}\quad La d\'emonstration donn\'ee dans la m\'emoire cit\'e est, on ne peut plus, simple et \'el\'egante. Elle expose d'abord sur un th\'eor\`eme connu (d\'emontr\'e au reste \S\ 6) qui peut \^etre \'enonc\'e ainsi; savoir, en supposant que les \'equations $U = 0$, $V = 0$, $\ldots$ soient satisfaites, on aura (pr\`es un facteur ind\'ependant de $\xi$, $\eta$, $\ldots$),
\[
\mathfrak{V}\phi = \xi\,(p\xi + q\eta + \cdots)^{\mu - \mu'}\bigl(\xi, \eta, \ldots\bigr),\quad
\mathfrak{V}\phi = \xi\,(p\xi + q\eta + \cdots)^{\mu - \mu''}\bigl(\xi, \eta, \ldots\bigr),\quad\&\mathrm{c}.
\]
Puis, cela fait, est fond\'ee sur la d\'emonstration d\'emontr\'ee (\S\ 12), savoir: le r\'esultant des fonctions
\begin{align*}
t\,(p\xi + q\eta + \cdots)^{\mu} + f'(\xi, \eta, \ldots),\\
t'(p\xi + q\eta + \cdots)^{\mu'} + f''(\xi, \eta, \ldots),\\
\vdots
\end{align*}
(o\`u $f$, $f'$, $\ldots$ sont des polyn\^omes de degr\'es $\mu$, $\mu'$, \ldots\ en $\xi$, $\eta$, \&c., des constantes quelconques, en supposant que $\mu$ ne soit plus petit qu'aucun autre des indices $\mu$, $\mu'$, \ldots, et en posant $\pi = \dfrac{\pi}{\mu} + \dfrac{\pi}{\mu'} + \cdots$, le r\'esultant $\Phi$ du syst\`eme sera tout au plus du degr\'e $\dfrac{\pi}{\mu}$ par rapport \`a $t$ et de m\^eme par rapport aux coefficients de $f$, $f'$, \ldots\ Le degr\'e par rapport \`a les coefficients est ai\,; le degr\'e par rapport aux coefficients de $f$, $f'$, \ldots\ sera donc au moins un infiniment petit de l'ordre $k$ contenu \`a $\phi^{\pi - \mu}$ comme \emph{facteur}. Dans cette supposition les \'equations $\mathfrak{V}\phi = 0$, $\mathfrak{V}\phi = 0$, \&c.\ deviendront:
\begin{align*}
t\,(p\xi + q\eta + \cdots)^{\mu} + f(\xi, \eta, \ldots) &= 0,\\
t'\,(p\xi + q\eta + \cdots)^{\mu'} + f'(\xi, \eta, \ldots) &= 0,\\
\vdots
\end{align*}
o\`u $f$, $f'$, \ldots\ sont des polyn\^omes de degr\'es $\mu$, $\mu'$, \ldots\ dont les coefficients sont infiniment petits par rapport \`a tous les autres plus du degr\'e $\dfrac{\pi}{\mu}$ par rapport \`a $t$ et de m\^eme par rapport aux coefficients de $f$, $f'$, \ldots\ Le degr\'e par rapport \`a ces coefficients est ai (assujettis \`a la seule condition ci-dessus mentionn\'ee) \'etant d'ailleurs arbitraires, on suit que le r\'esultant $\Phi$ contiendra en g\'en\'eral cette m\^eme puissance $\phi^{\pi - \mu}$ \emph{comme facteur}; ce qu'il s'agissait de d\'emontrer.

\medskip

\noindent (\emph{Le texte poursuit avec le d\'eveloppement de calculs explicites et la d\'emonstration compl\`ete pour des syst\`emes binaires et ternaires. La derni\`ere portion du texte (p.\ 184) contient quelques d\'etails sur les hautes puissances de $\phi^{\pi - \mu}$ qui apparaissent comme facteurs dans le r\'esultant compl\'ement, le tout encadr\'e par des observations sur l'importance du th\'eor\`eme dans le contexte des r\'esultants multidimensionnels.})

\selectlanguage{english}

\end{document}
